% ============================================================
% Related Work — main paper keeps only the two CLOSEST lines:
%   §2.1 Agent behavioral watermarking (mechanism-adjacent)
%   §2.2 Long-term memory systems + provenance (problem-adjacent)
% Everything else (RAG/corpus, LLM text, crypto primitives,
% memory-ops-as-actions extras) → Appendix \ref{app:related-ext}.
% ============================================================
\section{Related Work}
\label{sec:related}

\todo{One-sentence lede: \memmark{} sits at the intersection of two
research lines — behavioral watermarking for LLM agents (which gives
us the mechanism) and long-term memory systems with provenance
(which gives us the threat model). Extended coverage of RAG /
corpus-level watermarking, LLM text watermarking, and cryptographic
provenance primitives is deferred to
Appendix~\ref{app:related-ext}.}

% ============================================================
\subsection{Behavioral Watermarking for LLM Agents}
\label{sec:related-behavioral}
% ============================================================
\todo{Cite \agentmark~\cite{agentmark}, Agent
Guide~\cite{agentguide}, ActHook~\cite{acthook},
AGENTWM~\cite{agentwm}. All four embed at the visible-action layer
(planning, tool call, trajectory) and require the trajectory at
verify time. \memmark{} embeds at the latent state-evolution
decision layer, which survives when the trajectory is gone.}

\todo{Relation to \agentmark: we reuse its binning sampler as a
black-box component inside \S\ref{sec:method-sampler}; the
contribution is not the sampler but (a) where to intercept inside a
third-party memory SDK, (b) how to construct $(\candset, \probdist)$
from open-vocabulary structured-JSON LLM output, and (c) how to
compose audits across a cascade of intercepted calls per memory
event.}

\todo{Adjacent line — memory ops as actions
(\textsc{Memory-as-Action}~\cite{mem-as-action},
A-MAC~\cite{a-mac}). These optimize \emph{which} memory op to take;
\memmark{} embeds attribution \emph{into} the op. Orthogonal but
the closest neighbor on the "memory writes are first-class
decisions" axis.}

% ============================================================
\subsection{Long-Term Memory Systems and Provenance}
\label{sec:related-memory}
% ============================================================
\todo{Operating environment: \amem~\cite{amem},
\graphiti~\cite{graphiti}, MemOS~\cite{memos},
MemMachine~\cite{memmachine}. Two are the backends we wrap
(\amem{} agentic notes, \graphiti{} temporal graph); the other two
are MemMark-compatible via the §\ref{sec:method-formal} adapter.}

\todo{Provenance-aware designs:
\tiermem~\cite{tiermem} ships explicit provenance chains and
immutable source anchoring; MemOS \texttt{MemCube} carries
versioning / origin / lifecycle fields. Both are
\emph{field-stamping} provenance — they encode "who wrote this"
into metadata that the writer controls. \memmark{}'s contrast is
exactly this: we treat the snapshot-holder as the writer (adversarial
storage), so any field-based provenance is forgeable. \tiermem{} /
MemOS therefore appear in \S\ref{sec:rq3} as the
\texttt{signed-metadata-only} ablation point.}

\todo{Threat-side motivation: A-MemGuard~\cite{amemguard} and
MemoryGraft~\cite{memorygraft} catalog memory-poisoning attacks on
LLM agents — these motivate our \S\ref{sec:rq4} attack model
(\texttt{poisoning}, \texttt{manual\_edits}). The Mnemonic
Sovereignty survey~\cite{mnemonicsovereignty} covers governance /
access control / poisoning defense but explicitly does not cover
behavioral watermarking on memory-evolve decisions, supporting our
positioning.}

\todo{Direct KG-backend baseline: \kgmark~\cite{kgmark} is dynamic
knowledge-graph watermarking and is the closest single-point
comparison on KG-structured memory. \S\ref{sec:rq1}–\S\ref{sec:rq4}
compare \memmark{} vs \kgmark{} on the \graphiti{} backend.
\kgmark{} is not applicable to \amem{} (notes network, not strict
KG).}

% ---------- positioning table (stays in main paper) ----------
% ============================================================
% Table: Related Work positioning matrix (README §10.6)
% ============================================================
\begin{table*}[t]
\centering
\small
\setlength{\tabcolsep}{4pt}
\renewcommand{\arraystretch}{1.15}
\newcommand{\yes}{\textcolor{ForestGreen!80!black}{\checkmark}}
\newcommand{\no}{\textcolor{red!70!black}{\ding{55}}}
\newcommand{\partial}{$\circ$}
\begin{tabular}{lcccccc}
\toprule
Dimension &
\makecell{Agent Behavioral\\(\agentmark{}, etc.)} &
\makecell{RAG / Corpus\\(RAG-WM, etc.)} &
\makecell{\kgmark{}\\(KG-only)} &
\makecell{LLM Text\\(KGW, etc.)} &
\makecell{Memory Systems\\(\tiermem, MemOS)} &
\textbf{\memmark} \\
\midrule
Behavioral (vs static)        & \yes & \no  & \yes & \no  & N/A & \yes \\
Embedded in memory layer      & \no  & \no  & \yes & \no  & N/A & \yes \\
In-record verifiable (no ext.\ log) & \no  & \no  & \no  & \no  & partial & \yes \\
Robust to memory lifecycle    & \no  & partial & partial & \no  & N/A & \yes \\
Cryptographic audit           & \no  & \no  & \no  & \no  & metadata only & \yes \\
Backend-invariant             & N/A  & partial & \no  & N/A  & N/A & \yes \\
\bottomrule
\end{tabular}
\caption{\textbf{Positioning matrix}. No prior line simultaneously
covers behavioral $\times$ memory-layer $\times$ in-record
verification $\times$ lifecycle robust $\times$ cryptographic audit
$\times$ backend-invariant. \kgmark{} and \tiermem{} are the two
closest neighbors; \memmark{} integrates the keyed-sampling stance
of one with the provenance-aware stance of the other across three
structurally different memory backends. \emph{partial} = supported
only under restrictive assumptions; see \S\ref{sec:related} per-cell
discussion.}
\label{tab:related-positioning}
\end{table*}

